\documentclass[UTF8,12pt]{ctexart}
\usepackage[a4paper,margin=24mm]{geometry}
\usepackage{amsmath,amssymb,bm,graphicx,adjustbox}
\newcommand{\fitmath}[1]{\begin{adjustbox}{max width=\linewidth}$\displaystyle #1$\end{adjustbox}}
\setlength{\parindent}{0pt}
\setlength{\parskip}{.7em}
\sloppy
\allowdisplaybreaks
\begin{document}
\section*{2005 年考研数学三 · 参考答案}
\subsection*{2005年真题参考答案}

\subsection*{一、填空题}

（1）\(\displaystyle 2\)。（2）\(\displaystyle xy=2\)。（3）\(\displaystyle 2e\,\mathrm dx+(e+2)\,\mathrm dy\)。（4）\(\displaystyle \dfrac{\displaystyle 1}{\displaystyle 2}\)。（5）\(\displaystyle \dfrac{\displaystyle 13}{\displaystyle 48}\)。（6）\(\displaystyle a=0.4,\allowbreak b=0.1\)。


\subsection*{二、选择题}

（7）B。（8）A。（9）D。（10）B。（11）C。（12）A。（13）D。（14）C。


\subsection*{三、解答题}

（15）\(\displaystyle \dfrac{\displaystyle 3}{\displaystyle 2}\)。（16）\(\displaystyle x^2\dfrac{\displaystyle \partial^2g}{\displaystyle \partial x^2}-y^2\dfrac{\displaystyle \partial^2g}{\displaystyle \partial y^2}=\dfrac{\displaystyle 2y}{\displaystyle x}f^{\prime}(\dfrac{\displaystyle y}{\displaystyle x})\)。（17）\(\displaystyle \dfrac{\displaystyle \pi}{\displaystyle 4}-\dfrac{\displaystyle 1}{\displaystyle 3}\)。


（18）\(\displaystyle S(x)=\begin{cases}\dfrac{\displaystyle 1}{\displaystyle 2x}\ln\dfrac{\displaystyle 1+x}{\displaystyle 1-x}-\dfrac{\displaystyle 1}{\displaystyle 1-x^2},\allowbreak &x\in(-1,\allowbreak 1)\setminus\{0\},\allowbreak \\0,\allowbreak &x=0.\end{cases}\)


（19）证明略。（可考虑 \(\displaystyle F(x)=\displaystyle \int_0^xg(t)f^{\prime}(t)\,\mathrm dt+\displaystyle \int_0^1f(t)g^{\prime}(t)\,\mathrm dt-f(x)g(1)\)，\(\displaystyle x\in[0,\allowbreak 1]\)，证明 \(\displaystyle F^{\prime}(x)\le0\)。）


（20）\(\displaystyle a=2,\allowbreak b=1,\allowbreak c=2\)。


（21）（I）\(\displaystyle P^{\mathrm T}DP=\begin{pmatrix}A&O\\O&B-C^{\mathrm T}A^{-1}C\end{pmatrix}\)。（II）矩阵 \(\displaystyle B-C^{\mathrm T}A^{-1}C\) 是正定矩阵。证明略。（\(\displaystyle \begin{pmatrix}A&O\\O&B-C^{\mathrm T}A^{-1}C\end{pmatrix}\) 是正定矩阵，其各阶顺序主子式均大于零，证明 \(\displaystyle B-C^{\mathrm T}A^{-1}C\) 的各阶顺序主子式均大于零。）


（22）（I）\(\displaystyle f_X(x)=\begin{cases}2x,\allowbreak &0<x<1,\allowbreak \\0,\allowbreak &\text{其他},\allowbreak \end{cases}\quad f_Y(y)=\begin{cases}1-\dfrac{\displaystyle y}{\displaystyle 2},\allowbreak &0<y<2,\allowbreak \\0,\allowbreak &\text{其他}.\end{cases}\)（II）\(\displaystyle f_Z(z)=\begin{cases}1-\dfrac{\displaystyle z}{\displaystyle 2},\allowbreak &0<z<2,\allowbreak \\0,\allowbreak &\text{其他}.\end{cases}\)（III）\(\displaystyle \dfrac{\displaystyle 3}{\displaystyle 4}\)。


（23）（I）\(\displaystyle D(Y_i)=\dfrac{\displaystyle n-1}{\displaystyle n}\sigma^2\)，\(\displaystyle i=1,\allowbreak 2,\allowbreak \ldots,\allowbreak n\)。（II）\(\displaystyle \operatorname{Cov}(Y_1,\allowbreak Y_n)=-\dfrac{\displaystyle 1}{\displaystyle n}\sigma^2\)。（III）\(\displaystyle c=\dfrac{\displaystyle n}{\displaystyle 2(n-2)}\)。

\end{document}
